\begin{abstract}
O.~Caramello's \q{bridges} turn one property of a topos, read on two different presentations of it, into a theorem relating the two presentations; in a talk at Mines Paris (2024) she proposed to let a machine combine a database of such presentations with a database of such readings, and so generate theorems. We run a small version of that generator by hand and describe it as a software pipeline: two registries as inputs (five bridges proved in Lean~4 in the library \lp{} of lean-pool; small categories whose presheaves are familiar structures), then join, specialise, unfold, rewrite, verify (in Lean and by a direct proof) and filter against the literature. We report only outputs that we did not find in print: (1) representations of a quiver satisfy the weak excluded middle $\neg p\vee\neg\neg p$ exactly when every vertex has at most one outgoing arrow; (2) actions of a finite monoid satisfy it exactly when the monoid has a single minimal right ideal, which holds for the transition monoid of every strongly connected synchronizing automaton; (3) sets with an endomorphism have exactly three subtopoi, the middle one being the sets with an automorphism; (4) the essential points of the actions of a finite monoid are counted by its regular $\mathcal D$-classes, which gives $n$, $n+1$, $n+1$, $n$, $n+1$ of them for $T_n$, $PT_n$, $I_n$, $O_n$, $M_n(\mathbb F_q)$; (5) a transfer of the Ore condition to corners $eMe$; (6) the weak excluded middle for a dozen combinatorial topoi. Known outputs are logged separately for calibration, and a last section discusses the limits of the approach. A French version accompanies this text.
\end{abstract}
